% Audit des ajouts à mod_stage : corrections du premier audit, questionnaire de demande de stage,
% SIRET et validation de la demande par le référent ou le responsable de la thématique.
% Compilation : latexmk -pdf audit-des-ajouts.tex
\documentclass[11pt,a4paper]{article}
\usepackage[french]{babel}
\usepackage[T1]{fontenc}
\usepackage{lmodern,microtype,booktabs,tabularx,xcolor,hyperref,enumitem}
\usepackage[margin=2.2cm]{geometry}
\usepackage{url}
\expandafter\def\expandafter\UrlBreaks\expandafter{\UrlBreaks\do\a\do\b\do\c\do\d\do\e\do\f\do\g\do\h\do\i\do\j\do\k\do\l\do\m\do\n\do\o\do\p\do\q\do\r\do\s\do\t\do\u\do\v\do\w\do\x\do\y\do\z}
\DeclareUrlCommand\code{\urlstyle{tt}}
\Urlmuskip=0mu plus 2mu
\newcommand{\grav}[1]{\textbf{\textcolor{#1}{\textbullet}}}
\newcommand{\maj}{\grav{orange}}
\newcommand{\mineure}{\grav{gray}}
\newcommand{\crit}{\grav{red}}
\setlist[itemize]{itemsep=2pt,topsep=3pt}
\newcommand{\ok}{\textcolor{teal}{\textbf{corrigé}}}
\newcommand{\partiel}{\textcolor{orange}{\textbf{en partie}}}
\newcommand{\non}{\textcolor{gray}{\textbf{non traité}}}
\title{Audit des ajouts à \texttt{mod\_stage}}
\author{Sébastien Lefebvre}
\date{7 novembre 2026}

\begin{document}
\maketitle

\section{Objet et méthode}

Ce rapport fait suite à l'audit du code (\code{audit-du-module.tex}). Il couvre tout ce qui a été
ajouté depuis le commit \code{631e02b} sur la branche \code{ccr-62a5efc4-2fregj} :

\begin{itemize}
  \item les corrections du premier audit (commits \code{9c64315} et \code{817c714}) ;
  \item une fonctionnalité nouvelle (commit \code{2a7da10}) : un questionnaire de demande de stage par
    thématique, le numéro SIRET de l'organisme d'accueil, et la validation de la demande par
    l'enseignant référent, le responsable de la thématique, ou les deux ;
  \item les corrections issues du présent audit (commits \code{aa62f30} et \code{11e399b}).
\end{itemize}

\noindent Après ces ajouts, \code{mod_stage} est en version~0.6.0 (\code{2026110700}) et
\code{mod_stagesynthesis} en version~0.2.2 (\code{2026110700}).

\begin{itemize}
  \item \textbf{Deux revues indépendantes}, conduites sans connaître les conclusions de l'auteur
    des changements. La première a porté sur la nouvelle fonctionnalité : accès, logique, formulaires, schéma,
    sauvegarde, RGPD, courriels et tests. La seconde a recherché les régressions introduites par les
    corrections et vérifié, assertion par assertion, que les tests nouveaux et anciens passeraient.
  \item \textbf{Contrôles automatiques} : \code{php -l} sur tous les fichiers modifiés et
    \code{phpcs} avec le standard \code{moodle}. Les deux sont propres après correction.
  \item \textbf{Limite} : PHPUnit et Behat n'ont pas pu être exécutés ici. Le montage d'un Moodle
    local, qui aurait téléchargé et exécuté du code externe, a été refusé par la politique de la
    session. Les tests ont donc été relus, pas joués : la CI de GitHub (PostgreSQL, MariaDB,
    Behat) doit les jouer avant toute fusion.
\end{itemize}

\noindent\textbf{Échelle} : celle du premier rapport. \crit~critique, \maj~majeure,
\mineure~mineure.

\section{Synthèse}

Les deux revues ont trouvé \textbf{un défaut critique, huit majeurs et treize mineurs}. Tous ont
été corrigés, sauf un point qui relève d'une décision fonctionnelle (section~\ref{sec:decision}).
Le défaut critique était une régression des corrections : la vérification des gabarits de
convention ne chargeait pas TCPDF, si bien que tout gabarit passait pour illisible et qu'aucune
convention ne pouvait plus être générée. Le test qui devait le couvrir chargeait lui-même la
bibliothèque, ce qui masquait le défaut.

\begin{center}
\small
\begin{tabularx}{\linewidth}{@{}l X l@{}}
\toprule
& \textbf{Constat} & \textbf{État} \\
\midrule
\crit & Gabarits PDF toujours refusés (TCPDF non chargé) : plus aucune convention générable & \ok \\
\maj & Demande bloquée quand le paramètre ou les responsables changent après une validation & \ok \\
\maj & Responsable de thématique désinscrit : toujours attendu, demandes bloquées, courriels envoyés & \ok \\
\maj & Responsable de thématique : voit et modifie l'adresse, le téléphone et la date de naissance de l'étudiant & \ok \\
\maj & Durée retenue corrigée par la DEVE non enregistrée & \ok \\
\maj & Thématique bornée dont seule la dernière année a une durée : requis nul, thématique « atteinte » & \ok \\
\maj & Page des thématiques : formulaires imbriqués, enregistrement en masse cassé, suppression sans confirmation & \ok \\
\maj & Deux tests d'événements cassés par les descriptions en anglais & \ok \\
\maj & Deux tests de la tâche d'invitation cassés ; étudiant introuvable relu sans fin & \ok \\
\bottomrule
\end{tabularx}
\end{center}

\section{Nouvelle fonctionnalité}

\subsection{Ce qui a été construit}

\begin{itemize}
  \item \textbf{Questionnaire de demande de stage.} C'est un quatrième type de liste
    d'évaluation (\code{request}), que la thématique choisit (\code{stage_theme.requestlistid})
    comme ses autres formulaires. L'étudiant y répond en même temps qu'il demande sa convention,
    que ce soit depuis la demande seule ou à l'enregistrement du stage. Les réponses obligatoires
    et les valeurs des notes et des choix sont contrôlées côté serveur. Les réponses sont
    présentées à ceux qui valident, à la DEVE (revue de la convention) et dans le détail du stage.
    Si l'étudiant change de thématique en corrigeant une demande refusée, les réponses à l'ancien
    questionnaire sont retirées.
  \item \textbf{SIRET.} Le champ est facultatif (\code{stage_convention_detail.siret}). Les espaces
    et les points sont acceptés à la saisie ; le numéro est stocké sur 14 chiffres. La clé de Luhn
    est vérifiée, avec l'exception des établissements de La Poste. Le SIRET figure dans la demande,
    dans la revue, dans le détail du stage et sur la première page de la convention PDF. Il est
    exporté et restauré avec l'export global.
  \item \textbf{Validation avant la DEVE.} Le paramètre \og Qui valide la demande\fg{}
    (\code{stage.conventionvalidator}) vaut \code{referent}, \code{themeteacher} ou \code{both}.
    Il n'agit que si la validation préalable est activée. Chaque validation est enregistrée à part
    (\code{conventionteachervalidatedby} et \code{conventionthemevalidatedby}, avec leur date).
    La demande passe à la DEVE quand toutes les validations exigées sont réunies ; un refus
    l'arrête aussitôt. Un responsable de thématique sans étudiant attribué voit les demandes de sa
    thématique sur la page de validation enseignante, et chaque validateur concerné reçoit un
    courriel. Une thématique sans responsable actif est validée par le référent.
  \item \textbf{Écritures conditionnelles.} Une validation n'est écrite que si la demande est
    toujours en attente, si le stage n'est pas annulé et si cette validation n'a pas encore été
    donnée. Le mécanisme est celui du test de concurrence (\code{stage_update_entry_if()}) : deux
    validateurs simultanés sont sérialisés, et un refus intercalé n'est pas effacé.
  \item \textbf{Schéma.} L'étape de mise à jour \code{2026110700} ajoute quatre champs et une clé.
    Ils sont pris en compte dans la sauvegarde et la restauration (remappage de l'utilisateur et
    de la liste), dans le fournisseur RGPD (métadonnées, contextes, export, dissociation) et dans
    les journaux (l'événement de validation précise à quel titre).
  \item \textbf{Tests.} Neuf tests PHPUnit (\code{request_validation_test.php}) et deux scénarios
    Behat (\code{request_validation.feature}). Le premier scénario suit une demande de bout en
    bout : l'étudiant répond, un SIRET invalide est refusé, le responsable valide, la DEVE reçoit
    la demande. Le second fait valider successivement le référent puis le responsable.
\end{itemize}

\subsection{Constats et corrections}

\begin{itemize}
  \item[\maj] \textbf{Demande bloquée.} En mode \og les deux\fg{}, le référent valide, puis la DEVE
    repasse le paramètre à \og référent\fg{} ou retire le responsable. Plus aucune validation
    n'était attendue, mais la demande restait en attente, sans que personne puisse agir.
    \emph{Correction} : \code{stage_release_pending_conventions()} transmet à la DEVE les demandes
    qui n'attendent plus rien. Elle est appelée à l'enregistrement du paramètre, au changement des
    responsables, et à l'ouverture des pages de validation et des conventions.
  \item[\maj] \textbf{Responsable désinscrit.} Sa ligne de responsabilité subsiste. Il restait
    exigé alors qu'il ne pouvait plus valider, et il était relancé par courriel. \emph{Correction} :
    seuls les responsables qui ont encore la capacité dans l'activité comptent
    (\code{stage_get_active_theme_teachers()}), pour la règle comme pour les courriels.
  \item[\maj] \textbf{Coordonnées de l'étudiant.} Le formulaire de validation montrait au
    responsable de thématique l'adresse, le téléphone et la date de naissance de l'étudiant, et les
    rendait modifiables. Le détail du stage les lui masque pourtant. \emph{Correction} : pour un
    validateur qui n'est pas le référent, ces champs, le choix du référent et la case \og convention
    papier\fg{} du référent sont absents du formulaire et conservés tels quels à l'enregistrement.
  \item[\mineure] Siège de La Poste refusé ; SIRET contrôlé seulement sur \og Valider\fg{} et jamais à
    l'import global ; bouton \og Valider la demande\fg{} proposé à un référent qui n'avait rien à
    valider ; validation du responsable absente de l'export global et du détail du stage ; réponses
    au questionnaire supprimées au transfert entre promotions ; libellés et courriels qui ne
    parlaient que du référent ; commentaire inexact dans la revue DEVE. \emph{Tous corrigés}.
\end{itemize}

\subsection{Décision à prendre}
\label{sec:decision}

En mode \og les deux\fg{}, un enseignant qui est à la fois le référent de l'étudiant et le
responsable de la thématique donne les deux validations d'un seul geste. La demande initiale
(\og l'enseignant responsable de la thématique et/ou l'enseignant référent\fg{}) ne tranche pas. Le
comportement retenu considère que la même personne peut valider à ses deux titres. Si deux personnes
distinctes sont voulues, il suffit de limiter à un rôle par utilisateur dans
\code{stage_convention_user_validations()}.

\section{Corrections du premier audit}

\subsection{Régressions trouvées et corrigées}

\begin{itemize}
  \item[\crit] \textbf{Gabarits refusés.} \code{stage_convention_template_readable()} instanciait
    FPDI, qui étend TCPDF, sans charger \code{lib/pdflib.php}. L'erreur était attrapée et la fonction
    répondait toujours \og illisible\fg{} : tout gabarit était refusé à l'envoi, et toute génération
    s'arrêtait. \emph{Correction} : chargement de la bibliothèque dans la fonction, et retrait du
    \code{require_once} qui masquait le défaut dans le test.
  \item[\maj] \textbf{Durée retenue.} L'édition DEVE d'un stage validé n'écrit plus que la liste
    fermée des champs du stage (correction C10). La durée retenue, posée sur l'objet, n'était plus
    enregistrée. \emph{Correction} : elle est écrite à part.
  \item[\maj] \textbf{Thématique bornée.} La durée requise d'une thématique bornée à plusieurs années
    est désormais la somme des durées des années où l'étudiant a fait ses stages (correction C4).
    Si ces années n'avaient pas de durée propre, ou s'il n'y avait aucun stage, le requis tombait à
    zéro : la thématique passait pour atteinte. \emph{Correction} : repli sur la plus forte durée
    définie sur la plage, avec un test du cas.
  \item[\maj] \textbf{Formulaires imbriqués.} Le bouton de suppression en POST (correction A6) avait
    été placé dans le tableau des thématiques, déjà enfermé dans le formulaire d'enregistrement en
    masse. \emph{Correction} : les formulaires de suppression sont placés après le tableau, et chaque
    bouton les désigne par l'attribut \code{form}.
  \item[\maj] \textbf{Tests cassés.} Les assertions sur le texte des descriptions d'événements, passées en anglais,
    ont été mises à jour. Le test de la tâche d'invitation inscrit désormais son étudiant au cours. La tâche
    marque en échec, au lieu de l'ignorer, l'invitation d'un étudiant introuvable ou supprimé.
  \item[\mineure] Chaîne d'aide \og Retirer cette question\fg{} supprimée à tort parmi les 40 chaînes
    orphelines, car elle est utilisée dynamiquement : rétablie. Noms de thématique doublement échappés dans les
    doublons probables : corrigé. Test de la synthèse qui ne prouvait pas le contrôle de visibilité :
    réécrit avec une restriction d'accès.
\end{itemize}

\subsection{Bilan par constat du premier audit}

\begin{center}
\small
\begin{tabularx}{\linewidth}{@{}l X l@{}}
\toprule
\textbf{Constats} & \textbf{Traitement} & \textbf{État} \\
\midrule
B1--B3, B5 & Nettoyage et échappement, comptes rendus, entités hors HTML, fichiers temporaires & \ok \\
B4 & Noms de fichiers sans \code{s()} (non exploitable) & \non \\
C1--C8, C11, C13, C14 & Circuit, calculs, fuseau, transfert, longueurs, contournement, redemande, invitations, relance, dispense & \ok \\
C9 & Événements à l'import StageVet & \non \\
C10 & Écriture des seuls champs du formulaire et événement d'édition ; durée déclarée préremplie par les jours cochés & \partiel \\
C12, C15 & Durée retenue non bornée ; version de \code{install.xml} corrigée, autres points non & \partiel \\
D1--D5, D8 & Imports, StageVet, gabarit et logos, temporaires, classeurs invalides & \ok \\
D6 & Durée non entière et lignes incomplètes signalées ; UTF-16, fuseau de \code{parse_date}, mémoire : non & \partiel \\
D7, D11, D12 & Prévisualisations, collisions du zip, requêtes par ligne & \non \\
D9 & SIRET et validation du responsable ajoutés à l'aller-retour ; référents et réponses toujours absents & \partiel \\
D10 & Logo vérifié ; nom d'établissement par défaut conservé (choix de l'école) & \partiel \\
E1--E6, E9 & Registre, exports, dissociation, restauration des responsables, décodage des liens & \ok \\
E7, E10 & Noms dans le paramétrage ; API \code{core_external}, \code{readonlysession} & \non \\
E8 & Tests RGPD, restauration et \code{get_my_stages} en partie ajoutés & \partiel \\
E11 & Langue du destinataire et comptes supprimés ; \code{db/messages.php} et corps en \code{PARAM_RAW} : non & \partiel \\
A1--A5, A7--A9 & Accès, jeton après évaluation, minimisation, énumération, lectures avant connexion & \ok \\
A6 & Suppressions et transfert en POST ; réinitialisation et contournement DEVE restent des liens avec \code{sesskey} & \partiel \\
F1, F3, F7 & Antivirus, publication 0.6.0, documentation & \ok \\
F2, F4, F6 & Tests ajoutés, chaînes orphelines et événements en anglais, \code{.gitignore} ; \og ~:~\fg{} concaténés, seuil de couverture : non & \partiel \\
F5 & Découpage de \code{locallib.php} & \non \\
\bottomrule
\end{tabularx}
\end{center}

\noindent Les points non traités sont des mineures de maintenabilité ou de confort. Aucun ne
présente de risque pour les données. Ils restent à planifier, dans l'ordre du premier rapport.

\section{Ce qui tient}

\begin{itemize}
  \item Le contrôle d'accès de la nouvelle fonctionnalité. Un responsable d'une autre thématique
    est refusé, tout comme un référent quand seul le responsable est exigé. Une même validation ne
    peut pas être donnée deux fois, et une demande annulée ou refusée ne se valide plus.
  \item La concurrence : validations sérialisées, sans écrasement d'un refus.
  \item La concordance du schéma, de l'étape de mise à jour, de la sauvegarde, de la restauration
    et du fournisseur RGPD pour les quatre champs ajoutés.
  \item La majorité des corrections du premier audit, vérifiées une à une par la seconde revue,
    notamment l'échappement, le circuit, le fuseau des jours effectifs, l'antivirus, le RGPD et les
    accès.
  \item Les fichiers de langue : 850 clés de chaque côté, chacune définie en français et en anglais.
    Le code de style Moodle (\code{phpcs}) est respecté sans écart.
\end{itemize}

\section{Suite}

\begin{enumerate}[itemsep=2pt]
  \item Faire jouer la CI (PHPUnit sur PostgreSQL et MariaDB, Behat) sur la branche, et corriger ce
    qu'elle révélerait.
  \item Trancher la question de la section~\ref{sec:decision}.
  \item Planifier les mineures non traitées du premier audit.
\end{enumerate}

\end{document}
